\subsection{DEFINITIONS AND FIRST PROPERTIES}

Let G be a group and B a subgroup of G. The group \(B \times B\) acts on G by \((b, b')\) . \(g = bgb'^{-1}\) for \(b, b' \in B\) and \(g \in G\) . The orbits of \(B \times B\) on G are the sets BgB for \(g \in G\) , and are called the double cosets of G with respect to B. They form a partition of G; the corresponding quotient is denoted by B\textbackslash G/B. If C and C' are double cosets, CC' is a union of double cosets.

DEFINITION 1. A Tits system is a quadruple \((\mathbf{G},\mathbf{B},\mathbf{N},\mathbf{S})\) , where \(\mathbf{G}\) is a group, \(\mathbf{B}\) and \(\mathbf{N}\) are two subgroups of \(\mathbf{G}\) and \(\mathbf{S}\) is a subset of \(\mathrm{N} / (\mathrm{B}\cap \mathrm{N})\) , satisfying the following axioms:

(T1) The set \(\mathrm{B} \cup \mathrm{N}\) generates \(\mathrm{G}\) and \(\mathrm{B} \cap \mathrm{N}\) is a normal subgroup of \(\mathrm{N}\) .

(T2) The set S generates the group \(W = N/(B \cap N)\) and consists of elements of order 2.

(T3) \(sBw \subset BwB \cup BswB\) for \(s \in S\) and \(w \in W\) .

(T4) For all \(s \in S\) , \(sBs \not\subset B\) .

The group \(W = N / (B \cap N)\) is sometimes called the Weyl group of the Tits system \((G, B, N, S)\) .

Remarks. 1) We shall see in no. 5 (Cor. of Th. 3) that, if (G, B, N) is given, there exists at most one subset S of N/(B∩N) such that (G, B, N, S) is a Tits system.

\begin{enumerate}
\def\labelenumi{\arabic{enumi})}
\setcounter{enumi}{1}
\tightlist
\item
  Let \((G, B, N, S)\) be a Tits system, and let Z be a normal subgroup of G contained in B. Let \(G' = G/Z, B' = B/Z, N' = N/(Z \cap N)\) , and let \(S'\) be the image of S in \(N'/(B' \cap N')\) . Then one sees immediately that \((G', B', N', S')\) is a Tits system.
\end{enumerate}

Throughout this paragraph, with \((\mathrm{G},\mathrm{B},\mathrm{N},\mathrm{S})\) denoting a Tits system, we set \(T=B\cap N\) and W=N/T. A double coset means a double coset of G with respect to B. For any \(w\in W\) , we set \(C(w)=BwB\) ; this is a double coset.

We are going to deduce several elementary consequences of the axioms (T1) to (T4). We denote by \(w, w', \ldots\) elements of W and by \(s, s', \ldots\) elements of S. The following relations are clear:

\[
\mathrm{C} (1) = \mathrm{B}, \quad \mathrm{C} (w w ^ {\prime}) \subset \mathrm{C} (w). \mathrm{C} (w ^ {\prime}), \quad \mathrm{C} (w ^ {- 1}) = \mathrm{C} (w) ^ {- 1}.\tag{1}
\]

Axiom (T3) can also be written in the form

\[
\mathrm{C} (s). \mathrm{C} (w) \subset \mathrm{C} (w) \cup \mathrm{C} (s w).\tag{2}
\]

Moreover, since \(C(sw) \subset C(s).C(w)\) by (1) and since \(C(s).C(w)\) is a union of double cosets, there are only two possibilities:

\[
\mathrm{C} (s). \mathrm{C} (w) = \left\{ \begin{array}{l l} \mathrm{C} (s w) & \text {if} \mathrm{C} (w) \not \subset \mathrm{C} (s). \mathrm{C} (w) \\ \mathrm{C} (w) \cup \mathrm{C} (s w) & \text {if} \mathrm{C} (w) \subset \mathrm{C} (s). \mathrm{C} (w). \end{array} \right.\tag{3}
\]

By (T4), \(\mathrm{B} \neq \mathrm{C}(s). \mathrm{C}(s)\) ; putting \(w = s\) in (3) and using the relation \(s^2 = 1\) , we obtain

\[
\mathrm{C} (s). \mathrm{C} (s) = \mathrm{B} \cup \mathrm{C} (s).\tag{4}
\]

This formula shows that \(B \cup C(s)\) is a subgroup of G. Multiplying both sides of (4) on the right by \(C(w)\) , and using formula (3) and the relation

\[
\mathrm{B.C} (w) = \mathrm{C} (w),
\]

we obtain

\[
\mathrm{C} (s). \mathrm{C} (s). \mathrm{C} (w) = \mathrm{C} (w) \cup \mathrm{C} (s w).\tag{5}
\]

Taking the inverses of the sets entering into formulas (2), (3) and (5) and then replacing w by \(w^{-1}\) , we obtain the formulas

\[
\mathrm{C} (w). \mathrm{C} (s) \subset \mathrm{C} (w) \cup \mathrm{C} (w s)\tag{(2')}
\]

\[
\mathrm{C} (w). \mathrm{C} (s) = \mathrm{C} (w s) \quad \text {   if   } \mathrm{C} (w) \not \subset \mathrm{C} (w). \mathrm{C} (s)
\]

\[
\mathrm{C} (w). \mathrm{C} (s) = \left\{ \begin{array}{l l} \mathrm{C} (w) \cup \mathrm{C} (w s) & \text {if} \mathrm{C} (w) \subset \mathrm{C} (w). \mathrm{C} (s) \end{array} \right.\tag{3'}
\]

\[
\mathrm{C} (w). \mathrm{C} (s). \mathrm{C} (s) = \mathrm{C} (w) \cup \mathrm{C} (w s).\tag{(5')}
\]

Lemma 1. Let \(s_1, \ldots, s_q \in S\) and let \(w \in W\) . We have

\[
\mathrm{C} \left(s _ {1} \dots s _ {q}\right). \mathrm{C} (w) \subset \bigcup_ {\left(i _ {1}, \dots i _ {p}\right)} \mathrm{C} \left(s _ {i _ {1}} \dots s _ {i _ {p}} w\right),
\]

where \((i_1, \ldots, i_p)\) denotes the set of strictly increasing sequences of integers in the interval \([1, q]\) .

We argue by induction on \(q\) , the case \(q = 0\) being trivial. If \(q \geqslant 1\) , we have \(C(s_1 \ldots s_q).C(w) \subset C(s_1).C(s_2 \ldots s_q).C(w)\) . By the induction hypothesis, \(C(s_2 \ldots s_q).C(w)\) is contained in the union of the \(C(s_{j_1} \ldots s_{j_p}w)\) , where

\[
2 \leqslant j _ {1} <   \dots <   j _ {p} \leqslant q.
\]

By (T3), the set \(C(s_{1}).C(s_{j_{1}}\ldots s_{j_{p}}w)\) is contained in the union of the sets \(C(s_{1}s_{j_{1}}\ldots s_{j_{p}}w)\) and \(C(s_{j_{1}}\ldots s_{j_{p}}w)\) . This proves the lemma.

